% !TEX root = ../analisi-dei-capitolati.tex
% Capitolo 2 - Criteri e metodologia

\section{Criteri e metodologia}\label{sec:metodologia}

\subsection{Criteri di valutazione}
Per confrontare proposte molto diverse il gruppo ha usato sei criteri:
\begin{itemize}
  \item \textbf{interesse}: coinvolgimento dei membri verso il dominio e le
        tecnologie;
  \item \textbf{valore formativo}: tecnologie e metodi spendibili dopo il
        corso, in un contesto vicino a un prodotto reale;
  \item \textbf{chiarezza del perimetro}: precisione di requisiti, criteri
        di accettazione e prodotti attesi;
  \item \textbf{sostenibilità}: equilibrio tra sfida tecnica e carico di
        lavoro di sei persone nel budget orario del progetto;
  \item \textbf{vincoli esterni}: hardware da gestire, servizi a pagamento,
        tecnologie imposte, integrazioni non controllate dal gruppo;
  \item \textbf{supporto del proponente}: incontri periodici, canali
        rapidi, materiale di prova.
\end{itemize}

\subsection{Processo decisionale}
\begin{enumerate}
  \item \textbf{Lettura individuale} (12--13 ottobre): ogni membro ha letto
        i capitolati e approfondito quelli di suo interesse.
  \item \textbf{Prima selezione} (13 ottobre): dopo la presentazione e la
        discussione, i più votati sono stati C0, C1 e C2; una seconda
        votazione tra i soli finalisti ha dato l'ordine di preferenza della
        Tabella~\ref{tab:votazione}.
  \item \textbf{Approfondimento a coppie}: C0 a \nomi{cescobolla,corro}, C1
        a \nomi{biasini,guglielmo}, C2 a \nomi{mazzurana,spagnoli}. Ogni
        coppia ha preparato le domande per il proponente, validate dal
        gruppo.
  \item \textbf{Incontri con i proponenti}: Sanmarco Informatica il 21
        ottobre, sessione congiunta di Vimar il 22 ottobre, Eggon il 26
        ottobre.
  \item \textbf{Decisione finale}: conferma di C2, con C1 e C0 come
        alternative nell'ordine.
\end{enumerate}

\begin{table}[H]
  \centering
  \caption{Votazione tra i capitolati finalisti (13 ottobre 2026)\label{tab:votazione}}
  \begin{tabular}{@{}llrr@{}}
    \toprule
    \intestazione{Capitolato} & \intestazione{Proponente} &
    \intestazione{Voti} & \intestazione{Percentuale} \\
    \midrule
    \textbf{C2} -- NEXUM Concierge & Eggon S.r.l.                & 3 & 50\% \\
    \textbf{C1} -- AURA            & Vimar S.p.A.                & 2 & 33\% \\
    \textbf{C0} -- AIRA            & Sanmarco Informatica S.p.A. & 1 & 17\% \\
    \bottomrule
  \end{tabular}
\end{table}
